% ============================================
% Johns Hopkins Letter Template
% Assistant Professor Cover Letter – Middle Tennessee State University
% ============================================

\documentclass[11pt,letterpaper]{letter}

\usepackage{graphicx}
\usepackage{eso-pic}
\usepackage{geometry}
\usepackage{microtype}
\usepackage[hidelinks]{hyperref}

% ----- Page Setup -----
\geometry{left=25mm,right=25mm,top=25mm,bottom=25mm}

% ----- Logo Placement (first page only) -----
\newcommand{\PutJHULogoFG}[1]{%
  \AddToShipoutPictureFG*{%
    \AtPageUpperLeft{%
      \hspace{10mm}%
      \raisebox{-38mm}{%
        \includegraphics[width=6cm]{#1}%
      }%
    }%
  }%
}

% ----- Sender -----
\signature{Decory Edwards}
\address{Department of Economics \\ Johns Hopkins University \\ Baltimore, MD 21218 \\ \texttt{dedwar65@jh.edu}}

\begin{document}
\PutJHULogoFG{Images/jhulogo.png}

\begin{letter}{Search Committee \\
Department of Economics and Finance \\
Middle Tennessee State University \\
Murfreesboro, TN 37132 \\ USA}

\opening{Dear Members of the Search Committee,}

I am writing to express my interest in the tenure-track Assistant/Associate Professor position in Economics in the Department of Economics and Finance at Middle Tennessee State University, beginning August 1, 2026. I am a Ph.D. candidate in Economics at Johns Hopkins University, advised by Professors Christopher D.~Carroll, Nicholas W.~Papageorge, and Stelios Fourakis, and I expect to receive my degree in May 2026. My research fields are macroeconomics, wealth and income dynamics, and computational economics.

My research agenda focuses on understanding the determinants of wealth inequality and the mechanisms through which heterogeneity in returns to wealth shapes aggregate outcomes. In my job-market paper, I estimate the degree of heterogeneity in individual portfolio returns using micro-data from the Survey of Consumer Finances and simulate a structural model in which households face idiosyncratic return risk. I show that allowing for persistent heterogeneity in returns substantially improves the model’s ability to match the empirical distribution of wealth, offering a tractable way to connect micro-level return dispersion to macroeconomic inequality.

In a second project, I use the Health and Retirement Study to examine the relationship between interpersonal trust and portfolio performance. Building on the literature that finds a hump-shaped relationship between trust and income, I show that a similar relationship holds for individual returns to wealth measured in the HRS data, highlighting a behavioral channel through which social attitudes shape saving and investment outcomes. This work also provides new estimates of the persistent component of returns to net worth, which motivate the calibration of heterogeneity in my structural model.

MTSU’s commitment to innovative teaching, meaningful service, and robust research strongly resonates with my approach to academic work. I am excited by the opportunity to contribute to the department’s undergraduate and graduate programs by teaching across a range of fields, including \textit{Applied Microeconomics}, \textit{Macroeconomics}, \textit{Econometrics}, \textit{Labor Economics}, \textit{Industrial Organization}, and \textit{International Economics}. The position’s emphasis on applied and data-intensive research aligns closely with my training in computational macroeconomics and empirical microdata analysis, and I welcome the opportunity to incorporate modern tools such as Stata, Python, and R into coursework and research mentorship.

\textbf{Teaching.} My teaching experience centers on undergraduate macroeconomics and an upper-level macro policy seminar, where I have led weekly discussion sections and held regular office hours for classes ranging from 16 to 40 students. My approach is student-centered: I aim to help students “convince themselves” that they have moved from confusion to understanding by holding structured office-hour coaching that builds trust and confidence. At MTSU, I am prepared to teach undergraduate and graduate courses in macroeconomics, applied microeconomics, econometrics, and related elective fields, and I look forward to designing student-engaged, data-oriented learning experiences that reflect the department’s expectations.

I believe I would be a strong fit because my empirical and computational background allows me to (i) integrate modern data-analysis methods into economics instruction, (ii) supervise applied student research projects, and (iii) maintain a research agenda aimed at peer-reviewed publication in macroeconomics and empirical economics. I am especially drawn to MTSU’s commitment to serving a diverse student body and to fostering an environment where innovative teaching and inclusive mentorship are central to departmental life.

I would be honored to join the Department of Economics and Finance at Middle Tennessee State University. Thank you for your consideration; I welcome the opportunity to discuss my fit with the department at your convenience.

\closing{Sincerely,}

\end{letter}
\end{document}
